\documentclass{article}
\usepackage[T1]{fontenc}
\usepackage{lmodern}
\usepackage{graphicx}
\usepackage{amsmath,amssymb}
\usepackage[a4paper,margin=2cm]{geometry}
\usepackage[absolute,overlay]{textpos}
\setlength{\TPHorizModule}{1mm}
\setlength{\TPVertModule}{1mm}
\usepackage[indonesian]{babel}
\usepackage{hyperref}
\setlength{\emergencystretch}{2em}
% unit: Halaman 26

\begin{document}

% === Halaman 26 (llm) ===

\section*{RSA Encryption}

First of all we need the public key of the person to whom we want to send the message:

$(e, n) = (7, 77)$ (Enter appropriate values or Take public key from above)

Next we need the message. For simplicity let us demonstrate this here with just one letter. (In secure applications letters are never encrypted individually, but in whole blocks.)

Pick a letter to cipher: H

Before we can encrypt this letter we must encode it as a number:

$m = 8$ (here we take just the index in the alphabet)

Encryption itself is very simple: $m' = 57$ \[ (m' = m^e \pmod n) \]

The value $m'$ is the encrypted message sent to the receiver.

\vspace{10mm}

\section*{RSA Decryption}

First of all we need the private key of the person who got the encrypted message:

$(d, n) = (43, 77)$ (Enter appropriate values or Take private key from above)

Next we need the encrypted message:

$m' = 57$ (Enter an appropriate value or Take encrypted message from above)

Decryption again is very simple: $m = 8$ \[ (m = m'^d \pmod n) \]

The decoded message should match with the above chosen letter: H

\end{document}
